\section*{Model and main theorem (original Section 1)}
\begin{definition}[Algebraic Branching Programs]\label{def:abp}
An Algebraic Branching Program (ABP) is a layered graph where each edge is labeled by an affine linear form and the first and the last layer have one vertex each, called the ``start'' and the ``end'' vertex respectively.
The polynomial computed by an ABP is equal to the sum of the weights of all paths from the start vertex to the end vertex in the ABP, where the weight of a path is equal to the product of the labels of all the edges on it.
The size of an ABP is the number of vertices in it.
\end{definition}
A more general definition, which we sometimes adopt, is to allow every edge to be labeled by an arbitrary polynomial of degree at most $\Delta$. In this case we refer to the model as an ABP with edge labels of degree at most $\Delta$. Thus, the common case is $\Delta=1$, but our results are meaningful even when $\Delta=\omega(1)$.
\begin{theorem}\label{thm:intro main}
Let $\F$ be a field and $n\in\N$ such that $\Char(\F)\nmid n$. Then any algebraic branching program over $\F$ computing the polynomial $\sum_{i=1}^n x_i^n$ is of size at least $\Omega(n^2)$.
When the ABP's edge labels are allowed to be polynomials of degree at most $\Delta$, our lower bound is $\Omega(n^2/\Delta)$.
\end{theorem}
